\section{问题X：任务名称}

% 本模板是内容检查表，不要求所有问题拥有完全相同的子节；删除不适用项，禁止为齐全而填空话。
% ===== 一级问题开篇导读：生成正文时必须实例化，不能保留为注释或空壳 =====
% REPLACE_WITH_ACTUAL_QUESTION_GUIDE
% 用 2--4 句连续正文完成：
% 1. 回溯本问究竟在问什么、使用什么输入、是否承接前问；
% 2. 概括为什么采用当前路线，以及“先—再—最后”的方法顺序；
% 3. 说明将交付哪类结果/比较/决策及其后续用途。
% 不要照抄题面，不要写“本节将分析并得到结果”，不要在方法定义前堆公式或未锁定数字。

\subsection{任务定义与成功条件}

% 本节问题：本问究竟要计算、解释或决策什么，怎样才算答对？
% 输入/证据：题面条件、附件字段、约束编号、官方评价要求。
% 结果/接口：输出对象、单位、精度/可行性标准，以及与前后问的真实依赖。
用自己的话写清输入、已知条件、硬约束、目标输出和成功判据。说明与前后问题的真实依赖；后问不使用本问结果时，不得强造传递关系。

\subsection{数据审计与变量定义}

% 无数据附件或本问无需独立预处理时删除本小节。
% 本节问题：数据能否直接建模，哪些质量/单位/方向问题必须先处理？
% 输入/证据：样本清单、字段统计、缺失/异常/重复/时间空间范围、主键对齐结果。
% 结果/接口：清洗后的数据、变量字典、处理规则及其可能偏差。
报告样本量、字段类型、缺失/异常/重复、时间范围、单位、数据来源和跨表主键；给出变量含义、单位、取值域与方向。每项清洗、标准化或指标转向都说明原因，不能只写“进行了归一化”。

\subsection{基线与主模型}

% 本节问题：最简基线能做到什么，主模型额外解决了哪个真实缺口？
% 输入/证据：任务卡、机理/统计特征、基线与主模型同口径实测结果。
% 结果/边界：模型定义、假设、目标/损失、约束、参数来源，以及新增复杂度的代价。
先给最简可解释基线，再根据题面特征建立主模型。定义假设、变量、目标/损失、约束、参数来源和边界条件；未完成同口径实测的方案只能称候选，不能称改进。

\subsection{算法设计与求解}

% 本节问题：模型怎样被实际求出，别人按哪些设置可以复算？
% 输入/证据：算法/求解器配置、初始化、随机种子、停止条件、运行日志与复杂度记录。
% 结果/接口：可执行流程、关键参数、收敛/可行状态、运行成本和输出文件。
说明算法或求解器、初始化、参数/超参数、随机种子或重复次数、停止条件、约束处理、计算规模与运行成本。调用成熟工具时重点说明为什么适用和怎样配置，不粘贴无关源码。

\subsection{结果与题意解释}

% 本节问题：本问的直接答案是什么，证据相对基线怎样变化，为什么？
% 输入/证据：结构化结果、Figure/Table、配对比较或关键案例；所有数字回链结果文件。
% 结果/边界：答案、差值/区间、原因、负结果、代价与适用范围。
引用真实结构化结果，先直接回答本问，再按“数字/图表—与基线同口径比较—原因/机制—对模型或决策的含义—适用边界或代价”展开。每个 Figure 围绕一个科学问题组合互补 Panel，并在图后逐项解释关键面板及跨面板关系；不能只复述结果。图、表和文字中的数字必须回链到 \texttt{求解/问题X/结果/}，未经验证的原因只能标为待验证解释。

\subsection{验证与不确定性}

% 本节问题：核心结论经哪种独立检验仍成立，最敏感的风险是什么？
% 输入/证据：独立数据、残差/可行性、重复试验、灵敏度/消融/情景结果。
% 结果/边界：区间、失败场景、结论是否改变及触发换方案的条件。
按题型与风险选择独立互验、误差/残差、约束可行性、基线对比、灵敏度、鲁棒性、消融、收敛、蒙特卡洛或情景分析。明确数据隔离方式和不确定性来源；没有实验支撑时只报告待验证风险。

\subsection{本问结论与结果传递}

% 本节问题：用最短文字怎样直接回答题目，哪些结果可安全传给后问？
% 输入/证据：本问核心数字表、最终方案、适用范围和误差/不确定性。
% 接口：传递量、单位、文件路径、误差；没有真实依赖则删除“结果传递”。
给出可核验的本问答案和适用范围。若后续问题确实使用本问输出，列明传递量、单位、文件路径和不确定性；否则删除“结果传递”表述。
